\documentclass[12pt,a4paper]{article}

% ============================================================
% PACKAGES
% ============================================================
\usepackage[left=1in,right=1in,top=1in,bottom=1in]{geometry}
\usepackage{graphicx}
\usepackage{booktabs}
\usepackage{array}
\usepackage{tabularx}
\usepackage{colortbl}
\usepackage{xcolor}
\usepackage{makecell}
\usepackage{fancyhdr}
\usepackage{titlesec}
\usepackage{amsmath}
\usepackage{tcolorbox}

\definecolor{headerblue}{RGB}{213,220,228}

\newcolumntype{Y}{>{\centering\arraybackslash}X}
\newcolumntype{L}{>{\raggedright\arraybackslash}X}

\setlength{\parindent}{0pt}
\setlength{\parskip}{6pt}

\pagestyle{fancy}
\fancyhf{}
\rhead{SE: Machine Learning Lab}
\lhead{Lab 11}
\rfoot{\thepage}

\begin{document}

% ============================================================
% HEADER
% ============================================================
\renewcommand{\arraystretch}{2.0}
\begin{tabularx}{\textwidth}{|Y|}
\hline
\rowcolor{headerblue} \textbf{\large Department of Computer and Software Engineering} \\
\hline
\textbf{\large SE: Machine Learning} \\
\hline
\end{tabularx}
\renewcommand{\arraystretch}{1.0}

\vspace{12pt}

\renewcommand{\arraystretch}{2.0}
\begin{tabularx}{\textwidth}{|L|L|}
\hline
\textbf{Course Instructor:} & \textbf{Date:} \\
\hline
\textbf{Day:} & \textbf{Semester:} \\
\hline
\end{tabularx}
\renewcommand{\arraystretch}{1.0}

\vspace{12pt}

\begin{center}
{\Large\bfseries Lab 11: Refactoring, MLflow and Deployment}
\end{center}

\vspace{6pt}

\renewcommand{\arraystretch}{1.8}
\begin{tabularx}{\textwidth}{|Y|c|c|c|c|}
\hline
\textbf{Lab Title} & \textbf{CLO} & \textbf{BT} & \textbf{PLO} & \textbf{Weightage} \\
\hline
\makecell{{\small Refactoring \& Deployment} \\ {\small MLflow, APIs, Serialization}} & & & & \\
\hline
\end{tabularx}
\renewcommand{\arraystretch}{1.0}

\vspace{6pt}

\renewcommand{\arraystretch}{2.0}
\begin{tabularx}{\textwidth}{|Y|Y|Y|Y|Y|}
\hline
\textbf{Name} & \textbf{Roll Number} & \textbf{Report /10} & \textbf{Viva /5} & \textbf{Total /15} \\
\hline
 & & & & \\
\hline
\end{tabularx}
\renewcommand{\arraystretch}{1.0}

\vspace{12pt}

\begin{flushright}
Checked on: \_\_\_\_\_\_\_\_\_\_\_\_\_\_\_\_

\vspace{6pt}
Signature: \_\_\_\_\_\_\_\_\_\_\_\_\_\_\_\_
\end{flushright}

\newpage

\section*{Objectives}
\begin{itemize}
\item Refactor unstructured ML code into modular functions
\item Track experiments using MLflow
\item Serialize trained models
\item Build REST APIs using FastAPI
\end{itemize}

\section*{Problem Statement}

In real-world machine learning systems, models are not built as isolated scripts but as part of structured, maintainable, and deployable pipelines.

You are given a poorly written (``spaghetti'') machine learning script that mixes data loading, preprocessing, model training, and evaluation in a single block of code. Such code is difficult to maintain, extend, debug, and deploy.

Your task is to transform this code into a clean and modular pipeline by separating concerns into well-defined functions.

After refactoring, you will extend the pipeline to include:

\begin{itemize}
\item \textbf{Experiment Tracking:} Use MLflow to log model parameters, performance metrics, and artifacts.
\item \textbf{Model Serialization:} Save the trained model to disk and reload it for inference.
\item \textbf{Model Deployment:} Build a REST API using FastAPI to serve predictions from the trained model.
\end{itemize}

You will use the dataset from the previous lab (Lab 10) to ensure continuity.

By the end of this lab, you will have a complete mini machine learning system that includes:
\begin{itemize}
\item Clean, modular code
\item Experiment tracking
\item Persistent models
\item A working prediction API
\end{itemize}

\section*{Part A: Refactoring}

\begin{tcolorbox}
\textbf{Task 1: Identify Issues}

Open \texttt{train\_bad.py} and list at least four problems related to structure and maintainability:
\begin{enumerate}
    \item \textbf{Mixed Concerns:} I/O, preprocessing, and training are in the global scope, making deployment impossible without running the full script.
    \item \textbf{Magic Numbers:} Hardcoded parameters like \texttt{max\_depth=100} and \texttt{random\_state=42} are baked into the execution.
    \item \textbf{Data Leakage:} The target label was not properly isolated from the training features, leading to an artificially perfect accuracy.
    \item \textbf{Dead Code \& No Tracking:} Unused variables exist, categorical columns (\texttt{city}) are poorly handled, and no experiment metrics are preserved.
\end{enumerate}
\end{tcolorbox}

\vspace{6pt}

\begin{tcolorbox}
\textbf{Task 2: Modularization}

Refactored the script into isolated functions:
\begin{itemize}
  \item \texttt{load\_data()} — successfully handles purely I/O operations.
  \item \texttt{preprocess\_data(df)} — strictly drops the categorical \texttt{city} column and isolates \texttt{X} from \texttt{y} to prevent data leakage.
  \item \texttt{train\_model(X, y)} — fits the clean matrix without needing redundant data-type checks.
  \item \texttt{evaluate\_model(model, X, y)} — utilizes \texttt{sklearn.metrics} to return the float accuracy.
  \item \texttt{run\_pipeline()} — acts as the orchestrator to run the above pipeline in sequence.
\end{itemize}
\end{tcolorbox}

\section*{Part B: MLflow}

\begin{tcolorbox}
\textbf{Task 3: Logging}

Inside \texttt{run\_pipeline()}, the training phase is wrapped in an MLflow context manager:
\begin{itemize}
  \item \textbf{Parameters:} Logged \texttt{model\_type} and \texttt{random\_state} using \texttt{mlflow.log\_param}.
  \item \textbf{Metrics:} Logged the evaluation float via \texttt{mlflow.log\_metric("accuracy", acc)}.
  \item \textbf{Model:} Saved the model weights to the tracking server using \texttt{mlflow.sklearn.log\_model}.
\end{itemize}
\end{tcolorbox}

\section*{Part C: Serialization}

\begin{tcolorbox}
\textbf{Task 4: Save and Load Model}

Bridged the gap between the training script and the web API using \texttt{joblib}:
\begin{itemize}
  \item \textbf{Export:} Added \texttt{joblib.dump(model, "model.joblib")} at the end of the pipeline.
  \item \textbf{Import:} Added \texttt{joblib.load("model.joblib")} in the global scope of the FastAPI app, wrapped in a \texttt{try/except FileNotFoundError} block to prevent server crashes if booted before training.
\end{itemize}
\end{tcolorbox}

\section*{Part D: FastAPI}

\begin{tcolorbox}
\textbf{Task 5: API Endpoint}

Secured the \texttt{/predict} POST endpoint:
\begin{itemize}
  \item \textbf{Defense 1:} Raises a 503 error if \texttt{model is None}.
  \item \textbf{Defense 2:} Raises a 422 error if the \texttt{"features"} key is missing from the JSON.
  \item \textbf{Matrix Conversion:} Applied \texttt{np.array().reshape(1, -1)} to translate the 1D JSON array into a 2D matrix required by Scikit-Learn for inference.
\end{itemize}
\end{tcolorbox}

\section*{Discussion Questions}
\begin{itemize}
\item \textbf{Why is modular design important?} It prevents the "monolith trap", allowing individual components (like I/O or training) to be tested, scaled, and debugged independently without breaking the entire system.
\item \textbf{What are benefits of MLflow?} It acts as a version-control ledger for hyperparameter tuning, ensuring performance metrics and model artifacts are saved rather than lost when the RAM clears after a script finishes.
\item \textbf{Why is model serialization required?} It freezes the trained model structure into a binary file, allowing a web server (like FastAPI) to load the "brain" instantly without needing to retrain the model on every startup.
\end{itemize}

\end{document}